% ═══════════════════════════════════════════════════════════
% §6 Validation & Threats to Validity
% ═══════════════════════════════════════════════════════════
\FloatBarrier
\section{Validation \& Threats to Validity}
\label{sec:validation}

\subsection{Internal Validity}

The framework comparison (RQ1) rests on solid statistical ground: three independent seeds yield coefficients of variation below 2\% for both Unsloth ($\sigma$=2.6 \toksec{}) and HuggingFace ($\sigma$=0.8 \toksec{}), and the 5.2\% inter-framework difference produces Cohen's $d>3$, well beyond conventional significance thresholds. We are confident that the observed advantage is real, even if modest. The measurement-caliber correction (\S\ref{sec:statistics}) rescales all absolute values uniformly within each phase and therefore does not affect these relative conclusions.

The tuning sweep (RQ3) presents a more nuanced picture. Single-seed configurations with effect sizes exceeding 30\% are robust: the measurement variance established in B1 ($\sigma \approx 2\%$) provides a 15$\times$ safety margin against spurious findings. However, the interaction effect between G and sequence length (G=2 outperforming G=4 by only $\sim$4\% in the combined configuration) approaches the noise floor. We report this interaction as a hypothesis requiring multi-seed confirmation rather than an established fact.

Two methodological approximations warrant acknowledgment. The $6ND$ FLOPS model assumes all parameters participate equally in forward and backward computation; with LoRA training only 1.9\% of parameters, the backward-pass FLOPS are overestimated. This inflates absolute MFU values but does not affect \emph{relative} comparisons between configurations, which is our primary use case. Additionally, all measurements span 50--100 training steps; thermal throttling, memory fragmentation, and optimizer state growth that may emerge over thousands of steps are not captured.

\subsection{External Validity}

Our findings are scoped to a specific architectural intersection: RDNA3 (gfx1100), 3B-parameter models, and short-completion tasks. Each dimension of this scope carries transfer risk. RDNA4 introduces dedicated AI accelerators (2 per CU, native FP8) that fundamentally change the compute landscape; the dual-peak bias and quantization penalty we observe may not persist on that architecture. Scaling to 7B+ models alters GEMM aspect ratios and attention's memory footprint, potentially shifting the compute-bandwidth balance that underlies our gap attribution. Finally, GSM8K's short mathematical completions ($\leq$256 tokens) represent a favorable case for the sequence-length optimization; tasks requiring long-form generation (e.g., open-ended writing at 1024+ tokens) would see diminished returns from this lever.

The software stack (ROCm 7.2, Unsloth 2026.7.4, PyTorch 2.12.1) evolves rapidly. Absolute throughput values will shift with future releases. We argue, however, that the \emph{architectural insights} (bandwidth bottleneck dominance, quantization penalty mechanism, framework gain ceiling) and the \emph{methodology} (dual-peak MFU, component stacking, validity protocol) retain value independent of specific software versions.

\subsection{Construct Validity}

MFU measures hardware utilization, not learning efficiency. A configuration achieving 3.25\% \mfuprac{} is not necessarily producing a better-trained model than one at 2.3\%; the relationship between throughput and convergence speed (tokens-per-accuracy-point) is mediated by gradient quality, reward signal strength, and optimization dynamics that our profiler does not capture. We position this work as a \emph{systems efficiency} characterization: it answers "how fast can this hardware train?" rather than "how well does the resulting model perform?" Bridging these two questions through accuracy-aware efficiency metrics remains an important direction for future work.
